\documentclass[openany]{book}
\input{notes_white}
\usepackage{mathtools}
\usepackage{pgfplots}
\settitle{Proof of Concept - NeuroHTz}
\newcommand{\definitionbox}[1]{%
\par\noindent\fbox{%
\parbox{\dimexpr\linewidth-2\fboxsep-2\fboxrule\relax}{#1}%
}\par
}
\begin{document}
\title{\Large{\underemphasise{Proof of Concept - NeuroHTz}} \\ \vspace{5pt}
\large{Early Detection of Alzheimer's using Electroencephalography (EEG) and Machine Learning}}
\author{Eklavya Raman \\ \vspace{5pt} er24ms024@iiserkol.ac.in \\ 
\and Aman Mahnoori \\ \vspace{5pt} am24ms017@iiserkol.ac.in \\ 
\and Himanshu Mech \\ \vspace{5pt} hm24ms004@iiserkol.ac.in \\
\and Manjishta Debnath \\ \vspace{5pt} md24ms117@iiserkol.ac.in \\
\and Sourodeep Dey \\ \vspace{5pt} sd24ms011@iiserkol.ac.in \\
\and Akshat Mishra \\ \vspace{5pt} am24ms122@iiserkol.ac.in \\
\and Manideep Mondal \\ \vspace{5pt} mm24ms008@iiserkol.ac.in}
\date{\today}
\maketitle
\tableofcontents

\chapter*{Introduction}
    Alzheimer's disease is a progressive neurodegenerative disorder that primarily affects memory and cognitive function. Currently, accurate diagnosis of Alzheimer's disease is often made at later stages when significant brain damage has already occurred. Since there is no way to cure or halt the progression of the disease, early detection is crucial for managing symptoms and improving quality of life. If Alzheimer's can be detected at an early stage, patients can receive timely interventions, such as medications and lifestyle changes, that may slow down the progression of the disease. Additionally, early detection allows patients and their families to plan for the future and make informed decisions about care and support. To diagnose Alzheimer's disease, various methods are used, including clinical assessments, neuroimaging techniques, and biomarker analysis. However, these methods can be expensive, invasive, and may not always provide accurate results. \\ \\ The limitations of current diagnostic methods highlight the need for alternative approaches that are non-invasive, cost-effective, and accurate. In a country like India, where the healthcare infrastructure may be limited, especially in rural areas, there is a pressing need for accessible and affordable diagnostic tools for Alzheimer's and possibly other neurodegenerative diseases. Electroencephalography (EEG) is a non-invasive and cost-effective method that measures electrical activity in the brain. Recent research has shown that EEG can be used to detect early signs of Alzheimer's disease by analyzing patterns of brain activity. However, interpreting EEG data manually requires expertise and can be time-consuming. This is where machine learning comes into play. Machine learning algorithms can analyze large datasets of EEG recordings to identify patterns and features that are indicative of Alzheimer's disease. By training machine learning models on EEG data from patients with Alzheimer's and healthy controls, we can develop predictive models that can assist in the early detection of the disease. \\ \\
    We employ recent research in neuroscience, which provides proof that EEG can be used to detect early signs of Alzheimer's disease. Recent research also shows that 40Hz gamma oscillations are disrupted in Alzheimer's Disease, and that the steady state response power at 40Hz is a potential biomarker for early detection. We will leverage this and other relevant findings to guide our feature selection and model development process. Combining multiple approaches such as EEG data with and without 40Hz stimulation, and using different machine learning algorithms will allow us to develop a robust and accurate early detection tool for Alzheimer's disease. We also propose the use of cutting edge machine learning techniques and fuzzy logic to enhance the data pipeline, including artifact removal, feature extraction, and increase interpretability of the model. \\\\
    Since the main goal of this project is to develop a early detection tool for Alzheimer's disease that is accessible and affordable, our focus will be to use lightweight machine learning models that can be deployed on low-resource devices, such as smartphones or tablets. This will ensure that the tool can be used in a wide range of settings, including rural areas where access to healthcare may be limited. In this proof of concept, we will explore the feasibility of using EEG data and machine learning algorithms for the early detection of Alzheimer's disease. We will evaluate the performance of different machine learning models and identify key features in the EEG data that are associated with Alzheimer's disease. The ultimate goal is to develop a prototype tool that can be further refined and validated in clinical settings for widespread use in early detection of Alzheimer's disease, which would include both hardware and software components. We also aim to create a detailed, annotated dataset of EEG recordings from patients with Alzheimer's and other neurodegenerative diseases with and without 40Hz stimulation and make it open access, which can be used for future research and development in this field. By leveraging the power of machine learning and EEG data, we hope to contribute to the early detection and management of Alzheimer's disease, ultimately improving the lives of patients and their families.


\chapter{Proof of Concept}
\section{Abstract}
    In this proof of concept, we explore the feasibility of using electroencephalography (EEG) data and machine learning algorithms for the early detection of Alzheimer's disease. We leverage recent research that shows disruptions in 40Hz gamma oscillations in Alzheimer's patients and use this as a potential biomarker for early detection. We develop a data pipeline that includes artifact removal, feature extraction, and model development using lightweight machine learning algorithms suitable for deployment on low-resource devices. Our preliminary research indicates that EEG-based features can effectively differentiate between Alzheimer's patients and healthy controls, demonstrating the potential of this approach for early detection. We also propose a protocol for collecting and annotating EEG data from patients with Alzheimer's and other neurodegenerative diseases, which can be used for future research and development in this field. The ultimate goal of this proof of concept is to lay the groundwork for developing an accessible and affordable tool for early detection of Alzheimer's disease, particularly in resource-limited settings. \\
    Research shows the remarkable capability of machine learning algorithms to classify EEG data with high accuracy, making it a promising approach for early detection of Alzheimer's disease. The limitations of this, however, include the lack of longitudinal data and the need for large, annotated datasets to train machine learning models effectively. Additionally, while EEG is a non-invasive and cost-effective method, it can be susceptible to noise and artifacts, which can affect the accuracy of the models. We propose to address the limitation of longitudinal data by designing a protocol for collecting EEG data from patients at different stages of Alzheimer's disease, as well as from healthy controls. Using fuzzy logic, we can train machine learning models to handle the uncertainty and variability in EEG data, which can improve the robustness and interpretability of the models. Therefore, we might be able to develop a tool which is able to provide not only a binary classification of Alzheimer's vs. healthy control, but also a probabilistic assessment of the likelihood of Alzheimer's disease, which can be valuable for early detection and clinical decision-making. Potentially, this approach could also identify the stage of Alzheimer's disease in a patient, which is the main limitation with current studies due to the lack of longitudinal data. We will compare the performance of different models and feature sets to identify the most effective approach for early detection of Alzheimer's disease using EEG data.
\section{Methodology}
The NeuroHTz system follows a structured and multidisciplinary methodology that integrates neuroscience, biomedical hardware, signal processing, and machine learning to enable early-stage detection of Alzheimer's disease. The methodology is designed as a continuous pipeline, beginning with controlled neural stimulation and EEG acquisition, followed by signal preprocessing, feature extraction, machine learning analysis, and final risk assessment. \\
\subsection{Neural Stimulation and EEG Acquisition}
We conduct two types of EEG recordings: one with 40Hz auditory stimulation and one without any stimulation. Gamma oscillations, particularly at 40Hz, have been shown to be disrupted in Alzheimer's disease. Providing 40Hz auditory stimulation can enhance gamma oscillations in the brain, which may help to reveal differences in neural activity between Alzheimer's patients and healthy controls. This allows us to capture both the baseline neural activity and the brain's response to stimulation, providing a richer dataset for analysis. Further, the 40Hz data is shown to increase the accuracy of machine learning models in classifying Alzheimer's disease, as it provides a more distinct signal that can be used to differentiate between patients and controls. Comparing the results from both types of recordings will help us train models that can differentiate more effectively between Alzheimer's patients and healthy controls, both with and without stimulation. This dual approach allows us to leverage the potential benefits of 40Hz stimulation while also ensuring that our models can perform well even in the absence of stimulation, which is important for real-world applications where stimulation may not always be feasible. \\
Following stimulation, the system performs EEG signal acquisition using the custom-designed wearable headset. The hardware setup includes dry Ag/AgCl electrodes mounted on a headband, an analog front-end for amplification and filtering, and an ADS1299-based acquisition module for high-resolution digitization. The electrodes are strategically placed over frontal and temporal regions of the brain, which are most affected in the early stages of Alzheimer's disease. We shall use the International 10-20 system for electrode placement, ensuring consistency across recordings. The EEG signals are sampled at a rate of 500 Hz to capture the relevant frequency bands, including the gamma band around 40Hz. The analog front-end amplifies the microvolt-level EEG signals and applies essential filtering, including high-pass, low-pass, and notch filters, to isolate relevant brain activity and suppress noise. The digitized signals are then transmitted wirelessly via the ESP32 microcontroller to a processing unit for further analysis. We shall discuss the hardware design and specifications in detail in the subsequent sections, including the choice of components, circuit design, and data acquisition protocols.  \\
\subsection{Signal Preprocessing}
The raw EEG data acquired from the headset is often contaminated with various artifacts, such as eye blinks, muscle movements, and electrical noise. To ensure the quality of the data for analysis, we implement a comprehensive signal preprocessing pipeline. This includes bandpass filtering to isolate the relevant frequency bands (e.g., 1-100 Hz), notch filtering to remove power line noise (e.g., 50/60 Hz), and artifact removal techniques such as Independent Component Analysis (ICA) to separate and eliminate artifacts from the neural signals. We also explore the use of fuzzy logic-based approaches for artifact removal, which can help to handle the uncertainty and variability in EEG data more effectively and accurately. The preprocessed EEG data is then segmented into epochs corresponding to the stimulation periods and baseline periods for further analysis. This preprocessing step is crucial for improving the signal-to-noise ratio and ensuring that the features extracted from the EEG data are representative of the underlying neural activity related to Alzheimer's disease. We discuss the specific preprocessing techniques, parameters, and implementation details in the subsequent sections, along with the rationale for our choices and their impact on the quality of the data for machine learning analysis. \\
\subsection{Feature Extraction}
After preprocessing the EEG data, we extract relevant features that can be used for machine learning analysis. These features include time-domain features (e.g., mean, variance, skewness), frequency-domain features (e.g., power spectral density, band power in specific frequency bands), and time-frequency features (e.g., wavelet coefficients). We also focus on extracting features related to gamma oscillations, particularly around the 40Hz frequency, as disruptions in this band have been associated with Alzheimer's disease. Additionally, we explore the use of connectivity measures (e.g., coherence, phase-locking value) to capture the interactions between different brain regions, which may provide further insights into the neural changes associated with Alzheimer's disease. The extracted features are then normalized and organized into a feature matrix that serves as input for the machine learning models. We discuss the specific features extracted, their relevance to Alzheimer's disease, and the methods used for feature extraction in detail in the subsequent sections. We also explore the use of fuzzy features such as the Katz Fractal Dimension among others, which can capture the complexity of EEG signals and may provide additional discriminatory power for classifying Alzheimer's. \\
We further convert the raw signal data into a more interpretable and compact form using dimensionality reduction techniques such as Principal Component Analysis (PCA) or t-Distributed Stochastic Neighbor Embedding (t-SNE). This helps to reduce the dimensionality of the feature space while retaining the most important information, which can improve the performance of machine learning models and reduce the risk of overfitting. We also explore the use of fuzzy logic-based feature selection methods to identify the most relevant features for classification, which can enhance the interpretability and robustness of our models. \\
Finally, we convert raw EEG data into a visual representation using CWT (Continuous Wavelet Transform) scalograms, which can be used as input for convolutional neural networks (CNNs) and transformer based models. This hybrid approach allows us to leverage both traditional machine learning techniques on extracted features and deep learning techniques on raw data representations, providing a comprehensive analysis of the EEG data for early detection of Alzheimer's disease. We discuss the specific methods used for feature extraction, dimensionality reduction, and visual representation in detail in the subsequent sections, along with their rationale and impact on the performance of our machine learning models. \\
\subsection{Machine Learning Analysis}
We propose a hybrid approach for machine learning analysis that combines traditional machine learning algorithms with deep learning techniques. For the traditional machine learning approach, we explore algorithms such as Support Vector Machines (SVM), Random Forests, and Gradient Boosting Machines, which are known for their effectiveness in classification tasks. We train these models on the extracted features from the EEG data and evaluate their performance using metrics such as accuracy, precision, recall, and F1-score. We also implement cross-validation techniques to ensure the robustness of our models and to prevent overfitting. \\
For the deep learning approach, we explore the use of Convolutional Neural Networks (CNNs) and transformer-based models that can take the CWT scalograms as input. CNNs are particularly effective for image-based data and can capture spatial patterns in the scalograms, while transformer-based models can capture long-range dependencies in the data. We train these models using a suitable architecture and optimization techniques, and evaluate their performance using similar metrics as the traditional machine learning models. \\
We also explore the use of fuzzy logic-based machine learning models, such as fuzzy decision trees or fuzzy neural networks, which can handle the uncertainty and variability in EEG data more effectively. These models can provide not only a binary classification of Alzheimer's vs. healthy control but also a probabilistic assessment of the likelihood of Alzheimer's disease, which can be valuable for early detection and clinical decision-making. We compare the performance of different models and feature sets to identify the most effective approach for early detection of Alzheimer's disease using EEG data. We also discuss the interpretability of our models and how they can be used in clinical settings for early detection and risk assessment of Alzheimer's disease. \\
We propose the use of the Vision Transformer (ViT) architecture for the deep learning analysis of the CWT scalograms. The ViT architecture has shown remarkable performance in image classification tasks and can effectively capture spatial patterns in the scalograms that may be indicative of Alzheimer's disease. We will design a suitable ViT architecture, including the number of layers, attention heads, and embedding dimensions, and train it on our dataset of CWT scalograms. We will also implement data augmentation techniques to increase the diversity of our training data and improve the generalization of our model. The performance of the ViT architecture is experimentally demonstrated to be superior to traditional machine learning models and CNNs in classifying Alzheimer's disease from EEG data, particularly when using the 40Hz stimulation data, which provides a more distinct signal for classification. Since the ViT architecture using a global attention mechanism, it is superior to CNN's in capturing global patterns in the data, which can be crucial for identifying complex neural signatures associated with Alzheimer's disease. We will evaluate the performance of the ViT model using the same metrics as the traditional machine learning models and CNNs, and compare the results to identify the most effective approach for early detection of Alzheimer's disease using EEG data. \\
\subsection{Risk Assessment and Clinical Decision Support}
The final output of our machine learning analysis is a risk assessment for Alzheimer's disease, which can be used to support clinical decision-making. We propose to provide not only a binary classification of Alzheimer's vs. healthy control but also a probabilistic assessment of the likelihood of Alzheimer's disease, which can be valuable for early detection and clinical decision-making. This probabilistic assessment can help clinicians to identify patients who are at high risk of developing Alzheimer's disease and may benefit from further diagnostic testing or early interventions. We also explore the use of explainable AI techniques to enhance the interpretability of our models, allowing clinicians to understand the key features and patterns in the EEG data that are driving the predictions. We use empirical evidence from our machine learning analysis to identify the most important features and patterns in the EEG data that are associated with Alzheimer's disease, and we present these findings in a way that is accessible and informative for clinicians. We also propose the use of fuzzy logic-based decision support systems that can provide recommendations for further diagnostic testing or early interventions based on the risk assessment, which can enhance the clinical utility of our tool for early detection of Alzheimer's disease. We discuss the potential applications of our risk assessment and clinical decision support system in real-world clinical settings, particularly in resource-limited settings where access to healthcare may be limited. We also explore the ethical considerations and implications of using machine learning-based tools for early detection of Alzheimer's disease, including issues related to privacy, informed consent, and potential biases in the data and models. \\ 
We propose a user friendly interface for clinicians to interact with our risk assessment and clinical decision support system. This interface will allow clinicians to input patient data, view the results of the machine learning analysis, and receive recommendations for further diagnostic testing or early interventions based on the risk assessment. We will design the interface to be intuitive and accessible, ensuring that it can be easily used by clinicians with varying levels of technical expertise. We will also provide training and support materials to help clinicians understand how to use the tool effectively and interpret the results in a clinical context. The ultimate goal is to develop a tool that can be seamlessly integrated into clinical workflows for early detection of Alzheimer's disease, particularly in resource-limited settings where access to healthcare may be limited. \\
Since there is a lack of longitudinal data to support the ability of our proposed tool to provide definite diagnosis, it must be noted that the tool is intended to provide a risk assessment rather than a definitive diagnosis of Alzheimer's disease. The probabilistic assessment provided by our machine learning models can help clinicians to identify patients who are at high risk of developing Alzheimer's disease and may benefit from further diagnostic testing or early interventions. However, it is important to emphasize that the tool is not intended to replace clinical judgment or diagnostic testing, but rather to serve as a complementary tool that can assist clinicians in making informed decisions about patient care. We will clearly communicate the limitations of our tool and the importance of using it in conjunction with other clinical assessments and diagnostic tests for Alzheimer's disease. We also propose to conduct further research and validation studies to evaluate the performance of our tool in real-world clinical settings and to refine our models based on longitudinal data as it becomes available. \\

\section{Components and Implementation}


\chapter*{Appendix A - Features and Dimensionality Reduction}
\section{Feature Extraction and Selection}
We extract a comprehensive set of features from the preprocessed EEG data, including time-domain features (e.g., mean, variance, skewness), frequency-domain features (e.g., power spectral density, band power in specific frequency bands), and time-frequency features (e.g., wavelet coefficients). We also focus on extracting features related to gamma oscillations, particularly around the 40Hz frequency, as disruptions in this band have been associated with Alzheimer's disease. Additionally, we explore the use of connectivity measures (e.g., coherence, phase-locking value) to capture the interactions between different brain regions, which may provide further insights into the neural changes associated with Alzheimer's disease. We also explore the use of fuzzy features such as the Katz Fractal Dimension among others, which can capture the complexity of EEG signals and may provide additional discriminatory power for classifying Alzheimer's. \\
To understand this in more detail, we first define what features mean and why they are important. \\
\definitionbox{\textbf{Feature:} A feature is a measurable property or characteristic of data that can be used for analysis and modeling. In the context of EEG data, features can be derived from raw signals and can capture aspects of brain activity such as amplitude, frequency content, and connectivity patterns. Features are important because they serve as input to machine learning models, and the choice of features can significantly affect model performance. By extracting relevant features from EEG data, we provide machine learning models with information that helps differentiate between Alzheimer's patients and healthy controls.}
\definitionbox{\textbf{Feature Extraction:} Feature extraction is the process of transforming raw data into a set of features that can be used for analysis and modeling. In EEG analysis, feature extraction applies signal-processing techniques to derive information that is relevant to brain activity. For example, Fourier-based methods can be used for frequency-domain features such as power spectral density, and wavelet methods can be used for time-frequency features. The choice of technique depends on the characteristics of the EEG data and the goals of the analysis.}
\definitionbox{\textbf{Feature Selection:} Feature selection is the process of choosing a subset of relevant extracted features for use in machine learning models. This is important because too many features can lead to overfitting, where a model learns noise instead of underlying patterns. Feature selection improves performance by reducing dimensionality and focusing on informative variables. Common approaches include filter, wrapper, and embedded methods, and fuzzy logic-based selection can further improve interpretability and robustness under uncertainty.}
Let us denote the raw EEG data as a matrix organized by channels and time points. The feature extraction stage transforms this raw signal matrix into a sample-by-feature representation, and the feature selection stage keeps only the most informative subset for final classification. The resulting reduced feature set is then used as input to machine learning models. We discuss the specific extraction and selection techniques in the subsequent sections, along with their rationale and impact on early Alzheimer's detection using EEG data.
\section{Features Used in Alzheimer's Disease Detection}
We use a variety of features derived from the EEG data to capture different aspects of brain activity that may be relevant for Alzheimer's disease detection. The most standard features are - 
\begin{itemize}
    \item \textbf{Time-Domain Features:} These include statistical measures such as mean, variance, skewness, and kurtosis of the EEG signals. These features can capture the overall amplitude and variability of the brain activity.
    \item \textbf{Frequency-Domain Features:} These include power spectral density and band power in specific frequency bands (e.g., delta, theta, alpha, beta, gamma). Disruptions in specific frequency bands, particularly gamma oscillations around 40Hz, have been associated with Alzheimer's disease.
    \item \textbf{Time-Frequency Features:} These include wavelet coefficients that capture how the frequency content of the EEG signals changes over time. This can provide insights into the dynamic patterns of brain activity that may be relevant for Alzheimer's disease.
    \item \textbf{Connectivity Measures:} These include coherence and phase-locking value, which capture the interactions between different brain regions. Changes in connectivity patterns have been observed in Alzheimer's disease and may provide additional discriminatory power for classification.
    \item \textbf{Fuzzy Features:} These include measures such as the Katz Fractal Dimension, which can capture the complexity of EEG signals. Fuzzy features can provide additional information about the underlying neural dynamics that may be relevant for Alzheimer's disease detection.
\end{itemize}
\subsection{Time Domain Feature Extraction}
We extract time-domain features from the preprocessed EEG signals to capture the statistical properties of brain activity. These features include measures such as mean, variance, skewness, and kurtosis across channels and time segments. The mean reflects average signal amplitude, variance reflects variability, and skewness and kurtosis summarize distribution shape. These time-domain features help differentiate Alzheimer's patients from healthy controls by capturing differences in overall amplitude and variability patterns. We discuss the specific calculation methods and their relevance in the subsequent sections. The extraction process applies statistical operations to segmented EEG windows to produce a feature vector for each sample. We also explore fuzzy logic-based approaches for time-domain feature extraction to better handle uncertainty and variability while improving interpretability and robustness.
\subsection{Frequency Domain Feature Extraction}
We extract frequency-domain features from the preprocessed EEG signals to capture spectral properties of brain activity. These features include power spectral density and band power in standard EEG bands such as delta, theta, alpha, beta, and gamma. Disruptions in specific bands, particularly around 40 Hz in the gamma range, are associated with Alzheimer's disease. Spectral estimates can be obtained using methods such as FFT-based analysis or Welch averaging, and band-specific power values are aggregated for model input. These features help differentiate Alzheimer's patients from healthy controls by capturing changes in spectral structure. We discuss the specific methods and their relevance in the subsequent sections. We also explore fuzzy logic-based approaches for frequency-domain feature extraction to improve robustness and interpretability under noisy and variable EEG conditions.
\subsection{Time-Frequency Feature Extraction}
We extract time-frequency features from the preprocessed EEG signals to capture dynamic patterns of brain activity relevant to Alzheimer's disease detection. These features include wavelet coefficients obtained using Continuous Wavelet Transform analysis. This approach reveals how frequency content changes over time and provides temporal context that pure frequency summaries may miss. Different wavelet families, such as Morlet or Daubechies, can be used to emphasize different signal characteristics. The resulting coefficients are organized into structured feature matrices that combine time and frequency information for each sample and are then used as model inputs. We discuss the specific methods in later sections, and we also explore fuzzy logic-based variants to improve robustness and interpretability when EEG quality varies. \\
Further, we convert the raw signal data into a visual representation using CWT scalograms, which can be used as input for convolutional neural networks (CNNs) and transformer based models. This hybrid approach allows us to leverage both traditional machine learning techniques on extracted features and deep learning techniques on raw data representations, providing a comprehensive analysis of the EEG data for early detection of Alzheimer's disease. We will discuss the specific methods used for creating these visual representations and their relevance to our machine learning analysis in detail in the subsequent sections. The CWT scalograms provide a visual representation of the time-frequency characteristics of the EEG signals, which can be particularly useful for deep learning models that are designed to process image data, such as CNNs and Vision Transformers (ViTs). To do this, we can calculate the CWT coefficients as described above and then create a scalogram by plotting the magnitude of the coefficients as a function of time and frequency. This scalogram can be treated as an image input for our deep learning models, allowing them to learn complex patterns in the time-frequency domain that may be indicative of Alzheimer's disease. We will evaluate the performance of our deep learning models using these CWT scalograms and compare it to the performance of traditional machine learning models that use extracted features, to identify the most effective approach for early detection of Alzheimer's disease using EEG data.
\subsection{Connectivity Feature Extraction}
We extract connectivity features from the preprocessed EEG signals to capture interactions between brain regions that may reflect Alzheimer's-related neural changes. These features include coherence and phase-locking value, which summarize synchronization between channel pairs in complementary ways. Connectivity patterns are often altered in Alzheimer's disease and can therefore provide additional discriminatory signal for classification. We discuss the specific methods in subsequent sections. The extraction process applies connectivity analysis to segmented EEG windows and compiles channel-pair summaries into sample-level feature vectors for model input. We also explore fuzzy logic-based connectivity feature extraction to better handle uncertainty and variability while improving interpretability and robustness.

\section{Dimensionality Reduction}


\chapter*{Appendix B - Fuzzy Logic}
\section{Basic Overview of Fuzzy Logic}
Fuzzy logic is a mathematical framework for dealing with uncertainty and imprecision in data. Unlike traditional binary logic, which classifies statements as either true or false, fuzzy logic allows for degrees of truth, where a statement can be partially true and partially false at the same time. This is particularly useful in real-world applications where data may be noisy, incomplete, or ambiguous. In fuzzy logic, variables can take on a range of values between 0 and 1, representing the degree of membership in a particular set. For example, in the context of Alzheimer's disease detection, we might have a variable representing the likelihood of Alzheimer's disease, which can take on values such as "low," "medium," and "high" with corresponding membership functions that define how these linguistic terms relate to the underlying numerical values. Fuzzy logic systems typically consist of three main components: fuzzification, inference engine, and defuzzification. Fuzzification is the process of converting crisp input values into fuzzy values based on predefined membership functions. The inference engine applies a set of fuzzy rules to the fuzzified inputs to derive fuzzy outputs. Finally, defuzzification converts the fuzzy outputs back into crisp values that can be used for decision-making or further analysis. By using fuzzy logic in our machine learning models for Alzheimer's disease detection, we can enhance the interpretability and robustness of our models by allowing them to handle uncertainty and variability in EEG data more effectively. In this appendix, we provide a basic overview of fuzzy logic, including its principles, components, and applications in the context of our project for early detection of Alzheimer's disease using EEG data and machine learning algorithms. We also discuss how fuzzy logic can be integrated into our machine learning models to improve their performance and interpretability in the context of Alzheimer's disease detection.

\subsection{Basic priniciples and definitions}
In classical set theory, an element either belongs to a set or it does not. For example, in the case of a set of "Alzheimer's patients," an individual either belongs to the set (is diagnosed with Alzheimer's) or does not belong to the set (is not diagnosed with Alzheimer's). However, in many real-world scenarios, such binary classifications may not capture the complexity and uncertainty of the data. Fuzzy logic allows for a more nuanced representation of membership in a set by introducing the concept of fuzzy sets. In a fuzzy set, an element can have a degree of membership that ranges between 0 and 1. For example, an individual might have a membership value of 0.7 in the "Alzheimer's patients" set, indicating that they are likely to have Alzheimer's but there is some uncertainty. This allows us to capture the inherent variability and uncertainty in medical diagnoses and other real-world applications. Fuzzy logic also introduces the concept of linguistic variables, which are variables that can take on linguistic values such as "low," "medium," and "high." These linguistic values are associated with membership functions that define how they relate to the underlying numerical values. For example, we might have a linguistic variable for "risk of Alzheimer's disease" with membership functions that define what constitutes "low risk," "medium risk," and "high risk" based on the output of our machine learning models. By using fuzzy logic in our machine learning models for Alzheimer's disease detection, we can enhance their ability to handle uncertainty and variability in EEG data, which can improve their performance and interpretability in clinical settings. The formal definitions are as follows - \\
\definitionbox{\textbf{Fuzzy Set:} A fuzzy set in a universe of discourse is described by a membership function that assigns each element a degree of membership between zero and one. Unlike classical sets, where membership is strictly yes or no, fuzzy sets allow partial membership and therefore capture uncertainty more naturally.}
\definitionbox{\textbf{Linguistic Variable:} A linguistic variable is a variable whose values are words or labels such as low risk, medium risk, and high risk. Each label is represented by a fuzzy set and corresponding membership function that links language-based descriptions to numerical values.}
\subsection{Construction of Fuzzy Membership Functions}
The construction of fuzzy membership functions is a crucial step in the design of a fuzzy logic system. Membership functions define how the linguistic values of a variable relate to the underlying numerical values. There are various types of membership functions, including triangular, trapezoidal, Gaussian, and sigmoid functions. The choice of membership function depends on the specific application and the nature of the data. We can construct membership functions based purely on data-driven approaches, where we analyze the distribution of the data and define membership functions that capture the relevant patterns in the data. We can do this as follows - \\
\begin{enumerate}
    \item Collect a dataset of EEG recordings from patients with Alzheimer's disease and healthy controls, along with relevant clinical information such as age, gender, and cognitive test scores.
    \item Identify relevant features in the EEG data that are associated with Alzheimer's disease, such as power in specific frequency bands, connectivity measures, or complexity measures. 
    \item Use the identified features to define linguistic variables for the risk of Alzheimer's disease, such as "low risk," "medium risk," and "high risk."
    \item Use the minimize entropy priniciple approach to construct membership functions for the linguistic variables based on the distribution of the identified features in the dataset. This involves analyzing the distribution of the features and defining membership functions that capture the relevant patterns in the data while minimizing uncertainty and maximizing interpretability.
\end{enumerate}
\section{Fuzzy Features and Fuzzy Logic-Based Machine Learning Models}
Let us now discuss some fuzzy features that can be derived from EEG data and how they can be used in fuzzy logic-based machine learning models for Alzheimer's disease detection. A few examples of fuzzy features include -
\begin{itemize}
    \item \textbf{Katz Fractal Dimension:} This is a measure of the complexity of a signal, which can capture the irregularity and self-similarity of EEG signals. It can provide insights into the underlying neural dynamics and may be relevant for Alzheimer's disease detection.
    \item \textbf{Fuzzy Entropy:} This is a measure of the uncertainty or randomness in a signal, which can capture the variability and unpredictability of EEG signals. It can provide information about the complexity and irregularity of brain activity, which may be relevant for Alzheimer's disease detection.
    \item \textbf{Fuzzy Connectivity Measures:} These measures capture the interactions between different brain regions in a fuzzy manner, allowing for partial membership in connectivity patterns. This can provide insights into the changes in brain connectivity associated with Alzheimer's disease.
\end{itemize}
\end{document}